\chapter{Расширение: гибридный резервуарный слой для памяти Рин}
\label{ch:reservoirbranch}

\section{Почему эта ветка отличается от прочих будущих ветвей}

Не всякая дальняя ветвь одинаково близка к инженерной практике. Гиперграфы, клеточные
комплексы и квантовые носители полезны как сильные структурные и исследовательские идеи,
но резервуарно-динамическая ветвь выделяется на их фоне по трём причинам.

\begin{enumerate}
\item Она \textbf{многое объясняет} уже на языке текущей архитектуры: почему retrieval
можно понимать как возмущение системы, почему mood естественно трактовать как режим
динамики, почему replay похож на настоящую консолидацию, а не на summarization.
\item Она \textbf{интерпретируема}: здесь нет требования отдать весь смысл чёрному ящику.
Напротив, явный граф, быстрый индекс, динамический слой и write-path можно развести по ролям.
\item Она \textbf{doable}: первый прототип можно собрать, не переписывая весь мозг Рин,
а добавив рядом с текущим графовым позвоночником новый слой с чистым интерфейсом.
\end{enumerate}

\section{Ключевая идея: две формы памяти вместо одной}

Главный выигрыш этой ветки начинается не с замены графа, а с отказа от ложной альтернативы
``или архив, или динамика''. Плодотворная формула другая:
\[
\boxed{
\text{память Рин}=
\text{адресуемый архив}
\;+\;
\text{динамическая ассоциативная среда}
}
\]

Архив отвечает на вопросы:
\begin{itemize}
\item что действительно было записано;
\item когда это произошло;
\item откуда это известно;
\item что было удалено;
\item какие потомки надо ревизовать.
\end{itemize}

Динамическая среда отвечает на другие вопросы:
\begin{itemize}
\item что всплывает по частичному cue;
\item какие далёкие мосты сейчас собираются в одну конфигурацию;
\item насколько контекст ``живой'', а не просто найденный;
\item в какой ассоциативный бассейн скатывается текущее состояние.
\end{itemize}

\begin{principle}[Архив и динамика решают разные классы задач]
Если архив пытается быть живой памятью, он становится тяжёлой базой данных с плохой
ассоциативностью. Если динамика пытается быть архивом, она теряет provenance, deletion,
ревизию и юридическую адресность. Сильная архитектура не выбирает одно из двух, а
разводит их роли.
\end{principle}

\section{Краткая теория: что такое резервуар в этом контексте}

В классическом reservoir computing есть большая рекуррентная система с фиксированными
случайными связями:
\[
h_t = \sigma(W_{\mathrm{res}}h_{t-1}+W_{\mathrm{in}}x_t),
\]
и обучается только считывающий слой:
\[
y_t=W_{\mathrm{out}}h_t.
\]
Это даёт богатое нелинейное отображение истории входов в пространство состояний.

Но для памяти Рин классический резервуар в чистом виде недостаточен по двум причинам.
\begin{enumerate}
\item Он даёт в основном \emph{fading memory}, то есть естественную память на конечный
горизонт.
\item Он не решает задачу устойчивого хранения многолетнего опыта без изменения самого
хранилища.
\end{enumerate}

Поэтому здесь уместнее не ``reservoir computing как есть'', а гибрид:
\[
\boxed{
\text{быстрый ассоциативный индекс}
\;+\;
\text{медленный пластичный резервуар}
\;+\;
\text{явный архив}
}
\]
То есть:
\begin{itemize}
\item быстрый индекс обеспечивает pattern separation и запуск recall;
\item пластичный слой даёт богатую ассоциативную динамику;
\item архив страхует долгую память, provenance и точную ревизию.
\end{itemize}

\section{Архитектурный скелет}

\subsection*{1. Event encoder}
На вход идёт не ``голый факт'', а событие целиком: текст, время, настроение, источник,
социальный контекст, task-status, salience. Результатом становится вектор события $x_t$.

\subsection*{2. Быстрый разреженный индекс}
Из $x_t$ строится sparse code:
\[
z_t = \mathrm{kWTA}(Px_t),
\]
где активна малая доля компонент. Это ближайший аналог ``гиппокампального'' индекса:
не хранение всей памяти, а адресное зажигание стартовых площадок.

\subsection*{3. Быстрая ассоциативная связка}
На этом слое действует быстрая запись совместности:
\[
H \leftarrow H + \eta_h z_t z_t^\top,
\qquad \eta_h \text{ велика}.
\]
Он нужен для свежих связей, которые должны появляться почти немедленно.

\subsection*{4. Медленный пластичный резервуар}
Дальше живёт собственно динамический слой:
\[
r_{t+1}
=(1-\lambda)\,r_t + \lambda\,\sigma(W_t r_t + Ux_t + Bz_t).
\]
Пластичность медленная:
\[
W_{t+1}=W_t+\eta_c \Delta W_t,
\qquad \eta_c \ll \eta_h.
\]
Именно это разделение временных шкал делает идею похожей на CLS, а не на хаотическую
самопереписывающуюся сеть.

\subsection*{5. Replay}
В idle и sleep режимах архивные события снова проигрываются через индекс и резервуар.
Так пластичный слой учится не от разрозненных единичных вспышек, а от повторяемых,
структурно отобранных паттернов.

\subsection*{6. Readout}
Readout не должен быть ни тупым до бесполезности, ни настолько умным, чтобы съесть
всю работу динамики. Идеологически здесь есть три уровня:
\begin{enumerate}
\item простое считывание диагностик и классов;
\item считывание набора кандидатов / активных подграфов;
\item интерфейс в hidden-state или cross-attention слой LLM.
\end{enumerate}

\subsection*{7. Архив-учитель}
Явный граф остаётся не просто ``запасным вариантом'', а учителем и судьёй: он
поставляет материал для replay, хранит provenance и принимает окончательные записи.

\section{Новая инженерная рекомендация: начинать не со всего графа}

Здесь я специально добавляю предложение, которого не было в исходном обсуждении в такой
форме, но которое следует из него как наиболее разумный путь prototyping.

\begin{principle}[Первый резервуар должен жить на активном подграфе]
Самый практичный MVP — не ``пластичный океан поверх всей памяти'', а динамический слой
на активном подграфе, который уже поднят retrieval'ом: якоря, их окрестность, bridge-узлы,
несколько важных provenance-носителей и контрсвидетельства.
\end{principle}

Пусть retrieval поднял подграф $G_t=(V_t,E_t)$. Тогда резервуарный шаг идёт не на всём
$V$, а только на $V_t$:
\[
r^{(t)}_{k+1}
=(1-\lambda)\,r^{(t)}_k + \lambda\,\sigma(W^{(t)} r^{(t)}_k + u^{(t)}),
\]
где $W^{(t)}$ индуцирован текущим активным подграфом.

Преимущества этого решения:
\begin{enumerate}
\item вычислительная цена резко падает;
\item появляется хороший стенд для диагностики;
\item динамический слой автоматически привязан к интерпретируемому фрагменту памяти;
\item replay можно запускать локально по тематическим областям, а не по всему миру сразу.
\end{enumerate}

Недостаток только один: это уже не ``чистая распределённая память всего мозга'', а
слой локальной ассоциативной динамики. Но именно это и делает ветку инженерно подъёмной.

\section{Mood как контроллер режима}

Диалог очень точно попадает в одно важное место: mood в такой архитектуре полезно
понимать не как ещё один текстовый факт и не как косметический тег в промпте, а как
\emph{медленный регулятор режима динамики}.

Пусть $\gamma(m)$ — функция настроения. Тогда она может модифицировать эффективный
усилитель рекуррентной матрицы:
\[
W_{\mathrm{eff}} = \gamma(m)\,W.
\]
Или, более осторожно, влиять на leak, gain и bias перехода:
\[
r_{t+1}
=(1-\lambda(m))r_t
 + \lambda(m)\,\sigma\!\bigl(\beta(m)Wr_t + Ux_t + b(m)\bigr).
\]

\begin{intuition}[Почему именно mood подходит на эту роль]
У настроения есть свойство, которое в инженерии систем управления необычайно удобно:
оно инертно. Оно не дёргается на каждый токен, а меняется медленно. Значит, оно может
быть low-pass-контроллером, задающим общий ландшафт для быстрых когнитивных колебаний.
\end{intuition}

Интерпретация проста:
\begin{itemize}
\item спокойный режим делает динамику более диссипативной;
\item сфокусированный или тревожный режим подводит её ближе к зоне длинных ассоциаций;
\item творческий поиск может временно расширять допустимую область лёгкого хаоса.
\end{itemize}

\section{Гомеостатический контроллер: «пирожки Алисы»}

Вторая идея из диалога — ещё интереснее: не только медленный режимный сдвиг, но и
быстрый контур искусственного ``химического'' гомеостаза.

\subsection*{Градусник системы}
Полезно следить по крайней мере за двумя агрегатами:
\[
E(t)=\frac{1}{N}\sum_{i=1}^N r_i(t)^2,
\qquad
V(t)=\frac{1}{N}\sum_{i=1}^N (r_i(t)-r_i(t-1))^2.
\]
Первый измеряет ``энергию'' поля, второй — скорость его изменения.

\subsection*{Быстрые модуляторы}
Если система выходит из рабочего диапазона, можно временно впрыскивать глобальные
модуляторы.

\begin{enumerate}
\item \textbf{Седатив}: увеличивает утечку и уменьшает gain нелинейности.
\item \textbf{Кофеин}: уменьшает затухание и повышает возбудимость.
\item \textbf{Дофаминоподобный модулятор}: временно поднимает learning rate пластичности,
когда пришёл важный или неожиданный опыт.
\end{enumerate}

\begin{pitfall}[Это не подарок бесплатно]
Гомеостатический контур очень красив, но он вносит собственные параметры, собственные
границы, свои ошибки калибровки и свой новый способ сломать систему. Поэтому он не
должен появляться раньше, чем реально возникнет болезнь, которую он лечит.
\end{pitfall}

\section{Как объединяются mood и гомеостаз}

Лучший способ совместить эти две идеи — не путать их.

\begin{itemize}
\item \textbf{Mood} — медленный режимный слой. Он сдвигает целевые границы, в которых
система считает себя ``нормальной''.
\item \textbf{Гомеостаз} — быстрый автоматический контур. Он следит, чтобы текущая
динамика не вышла за границы, заданные более медленным режимом.
\end{itemize}

Например, в спокойном режиме $E_{\max}$ ниже и система быстрее ``успокаивается''. В
творческом режиме $E_{\max}$ выше: допускается более длинная и смелая ассоциативная
траектория, прежде чем включится торможение.

\section{Что строить сначала}

У этой ветви очень легко потерять масштаб и начать ``с будущего будущего''. Поэтому
полезен строгий порядок прототипирования.

\begin{enumerate}
\item \textbf{Шаг 1.} Пластичность на типизированном графе без отдельного резервуара:
Hebbian-like усиление со-активированных рёбер и replay от архива.
\item \textbf{Шаг 2.} Малый динамический слой на активном подграфе: проверка, даёт ли
он реальный выигрыш в pattern completion и bridge formation.
\item \textbf{Шаг 3.} Mood как параметр режима динамики.
\item \textbf{Шаг 4.} Только при реальной нестабильности — гомеостатические ``гормоны''.
\item \textbf{Шаг 5.} Более глубокий readout-интерфейс к LLM.
\end{enumerate}

\begin{workedbox}[Как выглядит хороший MVP]
Хороший MVP этой ветки — это не ``Рин теперь вся reservoir''. Хороший MVP — это
следующее: есть golden-микромир, есть текущий graph-retrieval, есть активный подграф,
есть малый пластичный слой поверх него, и можно строго измерить, улучшились ли
partial-cue recall, bridge rate и устойчивость к интерференции.
\end{workedbox}

\section{Как проверять}

Для этой ветки нужны отдельные тесты, отличные от обычного RAG-like retrieval.

\subsection*{1. Partial cue completion}
Подать не полный ключ, а его фрагмент, и измерить, дотягивается ли система до
правильной конфигурации лучше базовой версии.

\subsection*{2. Delay under interference}
Записать событие, потом зашумить систему похожими событиями, затем проверить, что
именно всплывает спустя время и после replay.

\subsection*{3. Provenance stress}
Проверить, не выросла ли способность вспоминать ценой потери объяснимости.

\subsection*{4. Stability envelope}
Смотреть на $E(t)$, $V(t)$, прокси для spectral radius и длину ассоциативных хвостов.

\subsection*{5. Ablation of controllers}
Сравнить: без mood-модуляции, с mood-модуляцией, с грубым homeostasis-контуром,
чтобы отделить полезную режимность от чисто декоративной.

\section{Главные риски}

\begin{enumerate}
\item \textbf{Readout bottleneck.} Динамика может быть богата, но недоступна для
хорошего извлечения.
\item \textbf{Пластичность против устойчивости.} Слишком живая система быстро учится и
быстро разрушается.
\item \textbf{Вычислительная цена.} Даже локальный резервуар нельзя бездумно крутить на
каждый запрос и каждый фоновой цикл.
\item \textbf{Конфликт с продуктовой адресностью.} Чем больше знания живёт в скрытой
динамике, тем тяжелее deletion, provenance и точная юридическая ревизия.
\end{enumerate}

\section{Вердикт}

Из всех ``будущих'' ветвей именно эта выглядит как самая содержательная промежуточная
ступень между сегодняшней графовой памятью и более радикальными субстратами.

\begin{principle}[Самая сильная версия ветки не радикальна, а гибридна]
Наиболее перспективная форма резервуарной ветви для Рин — это не отказ от графа, а
гибрид: явный архив, быстрый ассоциативный индекс, локально или тематически ограниченный
пластичный резервуар, replay и режимные контроллеры. Именно в этой форме она и
интерпретируема, и измерима, и инженерно правдоподобна.
\end{principle}

Если когда-нибудь эта ветвь действительно выстрелит, она даст Рин не ``ещё один
поисковый слой'', а новую форму живой памяти: менее буквальную, более ассоциативную,
но всё ещё связанную с архивом, который позволяет оставаться честной, редактируемой
и проверяемой.
